\section{Discussion and Limitations}
\label{sec:limitations}
The contribution is intentionally scoped. EARS surfaces multimodality; it does not create it.
The within-front sparsity factor can spread selection only across equivalent optima that already
exist. On weakly multimodal problems with essentially one good decision region, such as
shortest-path-style routing, the method has little to exploit. We therefore use the OpenStreetMap
facility-placement study as a real-map case study in the regime the method targets.

The opposite failure mode is also important. When decision-space diversity includes deceptive
\emph{local} optima, preserving diversity can be counterproductive. On the scalable MMF10--12
suite, EARS is not rank-1. It is fourth on the primary metric IGDX and third on the all-metric
average behind MMEA-WI and MO\_Ring (Section~\ref{sec:experiments}). This result defines the
method's main boundary: EARS is designed for equivalent global Pareto sets, not for deceptive
local basins.

The convergence-preservation claim is structural rather than asymptotic. The within-front
property (\S\ref{sec:method}) prevents a dominated solution from displacing a solution from a
better front, at a measured $3\%$ IGD cost in our ablation. It does not establish a global
convergence rate. The sparsity weight $\beta$ is also sensitive to problem geometry. Values
$\beta{\ge}1$ significantly regress IGD on MMF1/MMF5, so $\beta=0.5$ was frozen by
validation-only analysis (Figure~\ref{fig:beta}). This sensitivity affects the magnitude of the
coverage gain, not the placement claim itself. An adaptive $\beta$ driven by mode-occupancy
entropy is a natural extension.

The empirical scope is broad but not exhaustive. The benchmark suite covers MMF1--8, four CEC MMO
competition problems with verified analytic references, the deceptive MMF10--12 suite with
numerically validated global references, and scalable families up to $d{=}100$
(Section~\ref{sec:highdim}). MMF13--14 and the full IDMP family remain future work. The baselines
are faithful Python re-implementations validated on MMF, not the original MATLAB code, so small
fidelity gaps remain possible. Finally, the placement finding is specific to an unpenalised
sparsity signal under one Pareto-based backbone. The full behaviour of CPDEA and MMEA-WI also
reflects their operators, archives and convergence-penalised keys.

\section{Conclusions}
\label{sec:conclusion}
MMOP methods have often treated convergence and decision-space diversity as a trade-off to be
balanced. This work shows that the trade-off can be reduced when multimodality comes from
equivalent global Pareto sets. The key is to keep decision-space sparsity out of the dominance
sort and use it only as a non-negative intra-front selection term. Convergence is then carried by
rank, while decision-space diversity is carried by the within-rank factor.

A controlled experiment attributes the effect to this placement rather than to the multiplicative
form. The full EARS-MMOEA framework then turns the selection principle into competitive
performance on MMF benchmarks, a CEC MMO extension, scalable families and a five-city real-map OpenStreetMap
facility-placement case study. Ablations indicate that the equivalence-aware backbone is the main
driver, and real-map visualisations show EARS recovering more near-equivalent layouts than the
density baseline.

Within the tested scope, selection-layer separation provides a simple way to add decision-space
coverage at a small measured convergence cost. Future work should make $\beta$ adaptive, extend
the search-side mechanisms for weakly multimodal applications, and test larger benchmark suites
and more cities. All data, baselines, protocols, statistics and figures are released for full
reproducibility.
